\chapter{2010年四川卷（理）}

\section{单选题}

\begin{problem}
\(\i\) 是虚数单位，计算 \(\i + \i^2 + \i^3 =\)（\quad）

\choices
    {\(-1\)}
    {\(1\)}
    {\(-\i\)}
    {\(\i\)}
\end{problem}

\begin{answer}
A
\end{answer}

\begin{solution}
因为 \(\i^2=-1\)，所以 \(\i^3=-\i\). 因而
\[
\i+\i^2+\i^3=\i-1-\i=-1
\]
故选 A.
\end{solution}

\begin{problem}
下列四个图象所表示的函数，在点 \(x = 0\) 处连续的是（\quad）

\choices
    {\choicebitmap[width=0.15\paperwidth]{img/2010/sichuan_science/q2_optA_clean.png}}
    {\choicebitmap[width=0.15\paperwidth]{img/2010/sichuan_science/q2_optB_clean.png}}
    {\choicebitmap[width=0.15\paperwidth]{img/2010/sichuan_science/q2_optC_clean.png}}
    {\choicebitmap[width=0.15\paperwidth]{img/2010/sichuan_science/q2_optD_clean.png}}
\end{problem}

\begin{answer}
D
\end{answer}

\begin{solution}
函数在 \(x=0\) 处连续，需要满足在该点有定义，且当 \(x\to0\) 时函数值趋于该点的函数值. 图象 A 在 \(x=0\) 处没有定义，图象 B 在该点左右极限不相等，图象 C 虽然极限存在，但极限值与函数值不相等；只有图象 D 满足连续的定义.
故选 D.
\end{solution}

\begin{problem}
\(2\log_{5}10 + \log_{5}0.25 =\)（\quad）

\choices
    {\(0\)}
    {\(1\)}
    {\(2\)}
    {\(4\)}
\end{problem}

\begin{answer}
C
\end{answer}

\begin{solution}
由 \(0.25=\frac14\)，得
\[
2\log_{5}10+\log_{5}0.25
=\log_{5}100+\log_{5}\frac14
=\log_{5}25=2
\]
故选 C.
\end{solution}

\begin{problem}
函数 \( f(x) = x^2 + mx + 1 \) 的图象关于直线 \( x = 1 \) 对称的充要条件是（\quad）

\choices
    {\(m = -2\)}
    {\(m = 2\)}
    {\(m = -1\)}
    {\(m = 1\)}
\end{problem}

\begin{answer}
A
\end{answer}

\begin{solution}
二次函数 \(f(x)=x^2+mx+1\) 的图象的对称轴为
\[
x=-\frac{m}{2}
\]
题设要求对称轴为 \(x=1\)，所以 \(-\frac{m}{2}=1\)，即 \(m=-2\). 反之，当 \(m=-2\) 时对称轴确为 \(x=1\)，故这是充要条件.
故选 A.
\end{solution}

\begin{problem}
设点 \(M\) 是线段 \(BC\) 的中点，点 \(A\) 在直线 \(BC\) 外，\(\left|\overrightarrow{BC}\right|^{2} = 16\)，\(\left|\overrightarrow{AB} + \overrightarrow{AC}\right| = \left|\overrightarrow{AB} - \overrightarrow{AC}\right|\)，则 \(AM =\)（\quad）

\choices
    {\(8\)}
    {\(4\)}
    {\(2\)}
    {\(1\)}
\end{problem}

\begin{answer}
C
\end{answer}

\begin{solution}
因为 \(M\) 是 \(BC\) 的中点，所以
\(\overrightarrow{AB}+\overrightarrow{AC}=2\overrightarrow{AM}\)，\(\overrightarrow{AB}-\overrightarrow{AC}=\overrightarrow{CB}\).
由 \(\lvert\overrightarrow{BC}\rvert^2=16\) 得 \(\lvert\overrightarrow{CB}\rvert=4\). 于是
\[
2AM=\left\lvert2\overrightarrow{AM}\right\rvert
=\left\lvert\overrightarrow{AB}+\overrightarrow{AC}\right\rvert
=\left\lvert\overrightarrow{AB}-\overrightarrow{AC}\right\rvert=4
\]
所以 \(AM=2\).
故选 C.
\end{solution}

\begin{problem}
将函数 \( y = \sin x \) 的图象上所有的点向右平行移动 \( \frac{\pi}{10} \) 个单位长度，再把所得各点的横坐标伸长到原来的2倍（纵坐标不变），所得图象的函数解析式是（\quad）

\choices
    {\(y = \sin \left(2x - \frac{\pi}{10}\right)\)}
    {\( y = \sin \left(2x - \frac{\pi}{5}\right) \)}
    {\( y = \sin \left(\frac{1}{2} x - \frac{\pi}{10}\right) \)}
    {\( y = \sin \left(\frac{1}{2} x - \frac{\pi}{20}\right) \)}
\end{problem}

\begin{answer}
C
\end{answer}

\begin{solution}
向右平移 \(\frac{\pi}{10}\) 后，图象为
\[
y=\sin\left(x-\frac{\pi}{10}\right)
\]
再把横坐标伸长为原来的 \(2\) 倍，相当于把上式中的 \(x\) 换成 \(\frac{x}{2}\)，因此所得函数为
\[
y=\sin\left(\frac{x}{2}-\frac{\pi}{10}\right)
\]
故选 C.
\end{solution}

\begin{problem}
某加工厂用某原料由甲车间加工出 \(A\) 产品，由乙车间加工出 \(B\) 产品. 甲车间加工一箱原料需耗费 \(10\text{ 小时}\) 可加工出 \(7\text{ 千克}\) \(A\) 产品，每千克 \(A\) 产品获利 \(40\text{ 元}\). 乙车间加工一箱原料需耗费 \(6\text{ 小时}\) 可加工出 \(4\text{ 千克}\) \(B\) 产品，每千克 \(B\) 产品获利 \(50\text{ 元}\). 甲，乙两车间每天共能完成至多 \(70\text{ 箱}\) 原料的加工，每天甲，乙两车间耗费工时总和不得超过 \(480\text{ 小时}\)，甲，乙两车间每天总获利最大的生产计划为（\quad）

\choices
    {甲车间加工原料 \(10\text{ 箱}\)，乙车间加工原料 \(60\text{ 箱}\)}
    {甲车间加工原料 \(15\text{ 箱}\)，乙车间加工原料 \(55\text{ 箱}\)}
    {甲车间加工原料 \(18\text{ 箱}\)，乙车间加工原料 \(50\text{ 箱}\)}
    {甲车间加工原料 \(40\text{ 箱}\)，乙车间加工原料 \(30\text{ 箱}\)}
\end{problem}

\begin{answer}
B
\end{answer}

\begin{solution}
设甲、乙车间分别加工 \(x\)、\(y\) 箱原料，则
\(x\geq0\)，\(y\geq0\)，\(x+y\leq70\)，\(10x+6y\leq480\).
每加工一箱甲车间原料获利 \(7\times40=280\) 元，每加工一箱乙车间原料获利 \(4\times50=200\) 元，所以总利润为
\[
P=280x+200y
\]
可行域的顶点为 \((0,0)\)、\((0,70)\)、\((15,55)\)、\((48,0)\)，对应利润分别为 \(0\)、\(14000\)、\(15200\)、\(13440\) 元，最大值在 \((15,55)\) 处取得.
因此甲车间加工 \(15\) 箱、乙车间加工 \(55\) 箱.
故选 B.
\end{solution}

\begin{problem}
已知数列 \( \{a_{n}\} \) 的首项 \( a_{1} \neq 0 \)，其前 \(n\) 项和为 \( S_{n} \)，且 \( S_{n+1} = 2S_{n} + a_{1} \)，则 \( \displaystyle\lim_{n \to \infty} \frac{a_{n}}{S_{n}} = \)（\quad）

\choices
    {\(0\)}
    {\( \frac{1}{2} \)}
    {\(1\)}
    {\(2\)}
\end{problem}

\begin{answer}
B
\end{answer}

\begin{solution}
设 \(a_1=c\ne0\). 由 \(S_1=c\) 及递推式
\(S_{n+1}=2S_n+c\)，归纳得
\[
S_n=(2^n-1)c
\]
于是
\(a_n=S_n-S_{n-1}=2^{n-1}c\)，\(\frac{a_n}{S_n}=\frac{2^{n-1}}{2^n-1}\).
因此
\[
\lim_{n\to\infty}\frac{a_n}{S_n}=\frac12
\]
故选 B.
\end{solution}

\begin{problem}
椭圆 \(\frac{x^2}{a^2} + \frac{y^2}{b^2} = 1 \) （\(a > b > 0\)）的右焦点为 \(F\)，其右准线与 \(x\) 轴的交点为 \(A\)，在椭圆上存在点 \(P\) 满足线段 \(AP\) 的垂直平分线过点 \(F\)，则椭圆离心率的取值范围是（\quad）

\choices
    {\( \left(0,\frac{\sqrt{2}}{2}\right] \)}
    {\(\left(0, \frac{1}{2}\right]\)}
    {\(\left[\sqrt{2} - 1, 1\right)\)}
    {\( \left[\frac{1}{2},1\right) \)}
\end{problem}

\begin{answer}
D
\end{answer}

\begin{solution}
设椭圆焦距的一半为 \(c=\sqrt{a^2-b^2}\)，离心率为 \(e=\frac ca\)，则 \(0<e<1\)，右焦点为 \(F(c,0)\)，右准线为 \(x=\frac ae\)，从而 \(A\left(\frac ae,0\right)\).

点 \(F\) 在线段 \(AP\) 的垂直平分线上，等价于 \(FA=FP\). 对椭圆上点 \(P(x,y)\)，由焦点准线性质有
\[
FP=e\left(\frac ae-x\right)=a-ex
\]
所以
\(a-ex=FA=\frac ae-ae\)，\(\frac xa=1+\frac1e-\frac1{e^2}\).
椭圆上存在横坐标为 \(x\) 的点，当且仅当 \(-1\leq\frac xa\leq1\). 对 \(0<e<1\)，右侧不等式自动成立，而左侧给出
\[
1+\frac1e-\frac1{e^2}\geq-1
\Longleftrightarrow
2e^2+e-1\geq0
\Longleftrightarrow
e\geq\frac12
\]
故离心率的取值范围为 \(\left[\frac12,1\right)\).
故选 D.
\end{solution}

\begin{problem}
由 1，2，3，4，5，6 组成没有重复数字且 1，3 都不与 5 相邻的六位偶数的个数是（\quad）

\choices
    {\(72\)}
    {\(96\)}
    {\(108\)}
    {\(144\)}
\end{problem}

\begin{answer}
C
\end{answer}

\begin{solution}
六位偶数的末位只能是 \(2\)，\(4\)，\(6\)，所以不加限制时共有 \(3\times5!=360\) 个.

设事件 \(U\) 为数字 \(1,5\) 相邻，事件 \(V\) 为数字 \(3,5\) 相邻. 将相邻的两个数字看作一个整体，则
\[
|U|=|V|=3\times4!\times2=144
\]
若 \(U\)，\(V\) 同时发生，数字 \(1\)，\(5\)，\(3\) 必须组成 \(1\)，\(5\)，\(3\) 或 \(3\)，\(5\)，\(1\) 的整体，末位仍从 \(2\)，\(4\)，\(6\) 中选，故
\[
|U\cap V|=3\times3!\times2=36
\]
因此所求个数为
\[
360-144-144+36=108
\]
故选 C.
\end{solution}

\begin{problem}
半径为 \(R\) 的球 \(O\) 的直径 \(AB\) 垂直于平面 \(\alpha\)，垂足为 \(B\)，\(\triangle BCD\) 是平面 \(\alpha\) 内边长为 \(R\) 的正三角形，线段 \(AC\)，\(AD\) 分别与球面交于点 \(M\)，\(N\)，那么 \(M\)，\(N\) 两点间的球面距离是（\quad）

\bitmapfigure[max width=0.62\linewidth]{img/2010/sichuan_science/q11_fig1.png}

\choices
    {\(R\arccos \frac{17}{25}\)}
    {\(R\arccos \frac{18}{25}\)}
    {\(\frac{1}{3}\pi R\)}
    {\(\frac{4}{15}\pi R\)}
\end{problem}

\begin{answer}
A
\end{answer}

\begin{solution}
建立直角坐标系，取
\(B=(0,0,0)\)，\(A=(0,0,2R)\)，\(O=(0,0,R)\)
并令
\(C=(R,0,0)\)，\(D=\left(\frac R2,\frac{\sqrt3 R}{2},0\right)\).
直线 \(AC\) 上的点可写为 \(A+t(C-A)\). 代入球面方程
\(x^2+y^2+(z-R)^2=R^2\)，除去根 \(t=0\) 对应的点 \(A\) 外，另一根为 \(t=\frac45\). 因而
\[
\overrightarrow{OM}=\left(\frac{4R}{5},0,-\frac{3R}{5}\right)
\]
同理，直线 \(AD\) 与球面的另一交点满足
\[
\overrightarrow{ON}=\left(\frac{2R}{5},\frac{2\sqrt3 R}{5},-\frac{3R}{5}\right)
\]
两向量长度均为 \(R\)，故
\[
\cos\angle MON
=\frac{\overrightarrow{OM}\cdot\overrightarrow{ON}}{R^2}
=\frac{8R^2/25+9R^2/25}{R^2}
=\frac{17}{25}
\]
球面上两点间的短弧长等于半径乘以所对圆心角，所以球面距离为 \(R\arccos\frac{17}{25}\).
故选 A.
\end{solution}

\begin{problem}
设 \(a > b > c > 0\)，则 \(2a^2 + \frac{1}{ab} + \frac{1}{a(a - b)} - 10ac + 25c^2\) 的最小值是（\quad）

\choices
    {\(2\)}
    {\(4\)}
    {\(2\sqrt{5}\)}
    {\(5\)}
\end{problem}

\begin{answer}
B
\end{answer}

\begin{solution}
先合并倒数项：
\[
\frac1{ab}+\frac1{a(a-b)}=\frac1{b(a-b)}
\]
因为 \(a>b>0\)，有 \(b(a-b)\leq\frac{a^2}{4}\)，于是
\[
\frac1{b(a-b)}\geq\frac4{a^2}
\]
因此原式不小于
\[
2a^2+\frac4{a^2}-10ac+25c^2
=25\left(c-\frac a5\right)^2+a^2+\frac4{a^2}
\geq25\left(c-\frac a5\right)^2+4\geq4
\]
当 \(a=\sqrt2\)、\(b=\frac a2\)、\(c=\frac a5\) 时满足 \(a>b>c>0\)，且上述各步均取等号. 所以最小值为 \(4\).
故选 B.
\end{solution}

\section{填空题}

\begin{problem}
\(\left(2 - \frac{1}{\sqrt[3]{x}}\right)^6\) 的展开式中的第四项是\fillinblank{}.
\end{problem}

\begin{answer}
\(-160x^{-1}\)
\end{answer}

\begin{solution}
二项式展开的第四项对应于从第二项中取出 \(3\) 个因子，因此第四项为
\[
\binom63 2^{6-3}\left(-\frac1{\sqrt[3]{x}}\right)^3
=20\times8\times(-x^{-1})=-160x^{-1}
\]
\end{solution}

\begin{problem}
直线 \(x - 2y + 5 = 0\) 与圆 \(x^{2} + y^{2} = 8\) 相交于 \(A\)，\(B\) 两点，则 \(|AB| =\fillinblank{}\).
\end{problem}

\begin{answer}
\(2\sqrt{3}\)
\end{answer}

\begin{solution}
圆心为原点，半径为 \(2\sqrt2\). 圆心到直线 \(x-2y+5=0\) 的距离为
\[
d=\frac{|5|}{\sqrt{1^2+(-2)^2}}=\sqrt5
\]
弦长为
\[
|AB|=2\sqrt{(2\sqrt2)^2-d^2}=2\sqrt{8-5}=2\sqrt3
\]
\end{solution}

\begin{problem}
如图，二面角 \(\alpha - l - \beta\) 的大小是 \(60^{\circ}\)，\(AB \subset \alpha\)，\(B \in l\)，\(AB\) 与 \(l\) 所成的角为 \(30^{\circ}\)，则 \(AB\) 与平面 \(\beta\) 所成角的正弦值是\fillinblank{}.

\bitmapfigure[max width=0.72\linewidth]{img/2010/sichuan_science/q15_fig1.png}
\end{problem}

\begin{answer}
\(\frac{\sqrt{3}}{4}\)
\end{answer}

\begin{solution}
设直线 \(AB\) 与平面 \(\beta\) 所成的角为 \(\theta\). 在过 \(B\) 且垂直于棱 \(l\) 的截面内，\(AB\) 与棱的夹角为 \(30^{\circ}\)，二面角的截面角为 \(60^{\circ}\). 因而 \(AB\) 在垂直于 \(l\) 的方向上的分量为 \(AB\sin30^{\circ}\)，该分量在平面 \(\beta\) 的法线方向上的投影为 \(AB\sin30^{\circ}\sin60^{\circ}\). 由线面角的定义，
\[
\sin\theta=\sin30^{\circ}\sin60^{\circ}=\frac12\cdot\frac{\sqrt3}{2}=\frac{\sqrt3}{4}
\]
\end{solution}

\begin{problem}
设 \(S\) 为复数集 \(\mathbf{C}\) 的非空子集. 若对任意 \(x\)，\(y \in S\)，都有 \(x + y\)，\(x - y\)，\(xy \in S\)，则称 \(S\) 为封闭集. 下列命题：
\begin{circlelist}
\item 集合 \(S = \{a + b\i \mid a, b\text{ 为整数}, \i\text{ 为虚数单位}\}\) 为封闭集；
\item 若 \(S\) 为封闭集，则一定有 \(0 \in S\)；
\item 封闭集一定是无限集；
\item 若 \(S\) 为封闭集，则满足 \(S \subseteq T \subseteq \mathbf{C}\) 的任意集合 \(T\) 也是封闭集.
\end{circlelist}
其中真命题是\fillinblank{}. （写出所有真命题的序号）
\end{problem}

\begin{answer}
\(\circled{1}\)，\(\circled{2}\)
\end{answer}

\begin{solution}
命题①中，任取 \(a+b\i\)，\(c+d\i\in S\)，有
\[
(a+b\i)+(c+d\i)=(a+c)+(b+d)\i
\]
\[
(a+b\i)-(c+d\i)=(a-c)+(b-d)\i
\]
\[
(a+b\i)(c+d\i)=(ac-bd)+(ad+bc)\i
\]
由于整数的和、差、积仍为整数，三式都属于 \(S\)，所以命题①正确.

命题②中，取非空集合 \(S\) 中任意元素 \(x\). 由封闭性，\(x-x=0\in S\)，所以命题②正确.

命题③不正确，因为 \(S=\{0\}\) 是封闭集，但它是有限集.

命题④不正确. 取封闭集 \(S=\{0\}\)，再取 \(T=\{0,1\}\)，虽有 \(S\subseteq T\subseteq\mathbf C\)，但 \(1+1=2\notin T\)，所以 \(T\) 不是封闭集.

故真命题的序号为 \(\circled{1}\)，\(\circled{2}\).
\end{solution}

\section{解答题}

\begin{problem}
某种有奖销售的饮料，瓶盖内印有“奖励一瓶”或“谢谢购买”字样，购买一瓶若其瓶盖内印有“奖励一瓶”字样即为中奖，中奖概率为 \(\frac{1}{6}\)，甲，乙，丙三位同学每人购买了一瓶该饮料.
\begin{enumerate}
\item 求甲中奖且乙，丙都没有中奖的概率；
\item 求中奖人数 \( \xi \) 的分布列及数学期望 \( E_{\xi} \).
\end{enumerate}
\end{problem}

\begin{solution}
设每位同学中奖与否相互独立，中奖概率为 \(p=\frac16\)，未中奖概率为 \(q=\frac56\).

\begin{enumerate}
\item 甲中奖且乙、丙都未中奖的概率为
\[
p q^2=\frac16\left(\frac56\right)^2=\frac{25}{216}
\]
\item 中奖人数 \(\xi\) 可能取 \(0\)，\(1\)，\(2\)，\(3\). 由独立重复试验，
\[
P(\xi=0)=q^3=\frac{125}{216}
\]
\(P(\xi=1)=\binom31pq^2=\frac{75}{216}\)，\(P(\xi=2)=\binom32p^2q=\frac{15}{216}\)
\[
P(\xi=3)=p^3=\frac1{216}
\]
所以分布列为

\begin{center}
\begin{tabular}{c|cccc}
\(\xi\) & \(0\) & \(1\) & \(2\) & \(3\) \\
\hline
\(P\) & \(\frac{125}{216}\) & \(\frac{75}{216}\) & \(\frac{15}{216}\) & \(\frac{1}{216}\)
\end{tabular}
\end{center}

数学期望为
\[
E(\xi)=0\cdot\frac{125}{216}+1\cdot\frac{75}{216}+2\cdot\frac{15}{216}+3\cdot\frac1{216}=\frac12
\]
\end{enumerate}
\end{solution}

\begin{problem}
已知正方体 \(ABCD - A'B'C'D'\) 的棱长为 \(1\)，点 \(M\) 是棱 \(AA'\) 的中点，点 \(O\) 是对角线 \(BD'\) 的中点.
\begin{enumerate}
\item 求证：OM 为异面直线 \(AA'\) 和 \(BD'\) 的公垂线；
\item 求二面角 \( M - BC' - B' \) 的大小；
\item 求三棱锥 \(M - OBC\) 的体积.
\end{enumerate}

\bitmapfigure[max width=0.84\linewidth]{img/2010/sichuan_science/q18_fig1.png}
\end{problem}

\begin{solution}
\begin{enumerate}
\item 建立空间直角坐标系，取
\(A=(0,0,0)\)，\(\ B=(1,0,0)\)，\(\ C=(1,1,0)\)，\(\ D=(0,1,0)\)
\(A'=(0,0,1)\)，\(\ B'=(1,0,1)\)，\(\ C'=(1,1,1)\)，\(\ D'=(0,1,1)\).
于是
\(M=\left(0,0,\frac12\right)\)，\(O=\left(\frac12,\frac12,\frac12\right)\).
直线 \(AA'\) 的方向向量为 \((0,0,1)\)，直线 \(BD'\) 的方向向量为 \((-1,1,1)\)，而
\[
\overrightarrow{MO}=\left(\frac12,\frac12,0\right)
\]
因此
\(\overrightarrow{MO}\cdot(0,0,1)=0\)，\(\overrightarrow{MO}\cdot(-1,1,1)=0\).
又因为 \(M\in AA'\)，\(O\in BD'\)，所以 \(OM\) 同时垂直于两异面直线，且分别交于 \(M\)，\(O\)，故 \(OM\) 是它们的公垂线.

\item 二面角 \(M-BC'-B'\) 的棱为 \(BC'\). 平面 \(MBC'\) 内的两个向量可取
\(\overrightarrow{BC'}=(0,1,1)\)，\(qquad \overrightarrow{BM}=\left(-1,0,\frac12\right)\)
其法向量可取
\[
\mathbf n_1=\overrightarrow{BC'}\times\overrightarrow{BM}=\left(\frac12,-1,1\right)
\]
平面 \(BC'B'\) 为 \(x=1\)，法向量可取 \(\mathbf n_2=(1,0,0)\). 设二面角为 \(\theta\)，则
\[
\cos\theta=\frac{|\mathbf n_1\cdot\mathbf n_2|}{|\mathbf n_1|\,|\mathbf n_2|}
=\frac{1/2}{3/2}=\frac13
\]
由于 \(0<\theta<\frac\pi2\)，有
\[
\tan\theta=\frac{\sqrt{1-\cos^2\theta}}{\cos\theta}=2\sqrt2
\]
所以 \(\theta=\arctan(2\sqrt2)\).

\item 由
\(\overrightarrow{MB}=\left(1,0,-\frac12\right)\)，\(\overrightarrow{MC}=\left(1,1,-\frac12\right)\)，\(\overrightarrow{MO}=\left(\frac12,\frac12,0\right)\)
可得
\(\overrightarrow{MC}\times\overrightarrow{MO}=\left(\frac14,-\frac14,0\right)\)，\(\overrightarrow{MB}\cdot(\overrightarrow{MC}\times\overrightarrow{MO})=\frac14\).
因此三棱锥体积为
\[
V_{M-OBC}=\frac16\left|\overrightarrow{MB}\cdot(\overrightarrow{MC}\times\overrightarrow{MO})\right|=\frac1{24}
\]
\end{enumerate}
\end{solution}

\begin{problem}
\begin{enumerate}
\item
\begin{enumerate}
\item 证明两角和的余弦公式 \(C_{\alpha+\beta}\)：\(\cos (\alpha + \beta) = \cos \alpha \cos \beta - \sin \alpha \sin \beta\)；
\item 由 \(C_{\alpha+\beta}\) 推导两角和的正弦公式 \(S_{\alpha+\beta}\)：\(\sin (\alpha + \beta) = \sin \alpha \cos \beta + \cos \alpha \sin \beta\).
\end{enumerate}
\item 已知 \(\triangle ABC\) 的面积 \(S = \frac{1}{2}\)，\(\overrightarrow{AB} \cdot \overrightarrow{AC} = 3\)，且 \(\cos B = \frac{3}{5}\)，求 \(\cos C\).
\end{enumerate}
\end{problem}

\begin{solution}
\begin{enumerate}
\item
\begin{enumerate}
\item 在直角坐标系中，点 \(Q=(x,y)\) 绕原点逆时针旋转 \(\beta\) 后，其横坐标由两个正交投影相加得到，为 \(x\cos\beta-y\sin\beta\). 取单位圆上与正 \(x\) 轴夹角为 \(\alpha\) 的点，其坐标为 \((\cos\alpha,\sin\alpha)\). 旋转 \(\beta\) 后，该点的横坐标又等于 \(\cos(\alpha+\beta)\)，故
\[
\cos(\alpha+\beta)=\cos\alpha\cos\beta-\sin\alpha\sin\beta
\]
\item 利用 \(\sin t=\cos\left(\frac\pi2-t\right)\)，将 \(t=\alpha+\beta\) 代入上式，得
\[
\begin{gathered}
\sin(\alpha+\beta)=\cos\left(\frac\pi2-\alpha-\beta\right)\\
=\cos\left(\left(\frac\pi2-\alpha\right)+(-\beta)\right)\\
=\cos\left(\frac\pi2-\alpha\right)\cos(-\beta)-\sin\left(\frac\pi2-\alpha\right)\sin(-\beta)\\
=\sin\alpha\cos\beta+\cos\alpha\sin\beta.
\end{gathered}
\]
\end{enumerate}
\item 设三角形的内角为 \(A\)，\(B\)，\(C\)，并记 \(b=|AC|\)、\(c=|AB|\). 由面积和数量积条件，
\(bc\sin A=1\)，\(bc\cos A=3\).
因为 \(A\) 是三角形内角且 \(\cos A>0\)，所以
\(\sin A=\frac1{\sqrt{10}}\)，\(\cos A=\frac3{\sqrt{10}}\).
又由 \(\cos B=\frac35\) 及 \(0<B<\pi\)，得 \(\sin B=\frac45\). 因为 \(C=\pi-(A+B)\)，所以
\[
\cos C=-\cos(A+B)=\sin A\sin B-\cos A\cos B
=\frac4{5\sqrt{10}}-\frac9{5\sqrt{10}}
=-\frac1{\sqrt{10}}=-\frac{\sqrt{10}}{10}
\]
\end{enumerate}
\end{solution}

\begin{problem}
已知定点 \(A(-1,0)\)，\(F(2,0)\)，定直线 \(l: x = \frac{1}{2}\)，不在 \(x\) 轴上的动点 \(P\) 与点 \(F\) 的距离是它到直线 \(l\) 的距离的 \(2\) 倍. 设点 \(P\) 的轨迹为 \(E\)，过点 \(F\) 的直线交 \(E\) 于 \(B\)，\(C\) 两点，直线 \(AB\)，\(AC\) 分别交 \(l\) 于点 \(M\)，\(N\).
\begin{enumerate}
\item 求 \(E\) 的方程；
\item 试判断以线段 \(MN\) 为直径的圆是否过点 \(F\)，并说明理由.
\end{enumerate}
\end{problem}

\begin{solution}
\begin{enumerate}
\item 设 \(P=(x,y)\). 由题设
\[
PF=2\,d(P,l)
\]
其中 \(F=(2,0)\)，直线 \(l\) 为 \(x=\frac12\)，故
\[
(x-2)^2+y^2=4\left(x-\frac12\right)^2
\]
展开并整理得
\[
x^2-\frac{y^2}{3}=1
\]
所以轨迹 \(E\) 的方程为 \(x^2-\frac{y^2}{3}=1\)，并保留题设中的 \(y\ne0\) 条件.

\item 先考虑过 \(F\) 的非竖直直线，设其方程为 \(y=k(x-2)\). 由于它与 \(E\) 有两个交点，所讨论的情形不含 \(k=0\) 和 \(k^2=3\). 设两个交点 \(B\)，\(C\) 的横坐标为 \(x_1\)，\(x_2\). 代入轨迹方程得
\[
(3-k^2)x^2+4k^2x-(4k^2+3)=0
\]
从而
\(x_1+x_2=-\frac{4k^2}{3-k^2}\)，\(qquad x_1x_2=-\frac{4k^2+3}{3-k^2}\).
设 \(X_i=(x_i,k(x_i-2))\). 直线 \(AX_i\) 与 \(l\) 的交点横坐标为 \(\frac12\)，因此其纵坐标为
\(y_i=\frac{3k(x_i-2)}{2(x_i+1)}\)，\((i=1,2)\).
这里 \(y_1\)，\(y_2\) 分别是点 \(M\)，\(N\) 的纵坐标. 由韦达定理，
\[
\frac{(x_1-2)(x_2-2)}{(x_1+1)(x_2+1)}
=\frac{x_1x_2-2(x_1+x_2)+4}{x_1x_2+(x_1+x_2)+1}
=-\frac1{k^2}
\]
于是
\[
y_1y_2=\frac{9k^2}{4}\cdot\left(-\frac1{k^2}\right)=-\frac94
\]
因此
\(\overrightarrow{FM}=\left(-\frac32,y_1\right)\)，\(qquad \overrightarrow{FN}=\left(-\frac32,y_2\right)\)，\(qquad \overrightarrow{FM}\cdot\overrightarrow{FN}=\frac94+y_1y_2=0\).
若过 \(F\) 的直线为竖直直线 \(x=2\)，则其与 \(E\) 的交点为 \((2,3)\)、\((2,-3)\)，相应地 \(M\)，\(N\) 的纵坐标为 \(\frac32\)，\(-\frac32\)，同样有 \(\overrightarrow{FM}\cdot\overrightarrow{FN}=0\).

所以 \(\angle MFN=90^{\circ}\). 根据圆的直径所对的圆周角为直角，点 \(F\) 在以 \(MN\) 为直径的圆上.
\end{enumerate}
\end{solution}

\begin{problem}
已知数列 \(\{a_{n}\}\) 满足 \(a_1 = 0\)，\(a_2 = 2\)，且对任意 \(m\)，\(n \in \mathbf{N}^*\) 都有 \(a_{2m-1} + a_{2n-1} = 2a_{m+n-1} + 2(m - n)^2\).
\begin{enumerate}
\item 求 \(a_{3}\)，\(a_{5}\)；
\item 设 \( b_{n} = a_{2n + 1} - a_{2n - 1}\)（\(n \in \mathbf{N}^{*}\)），证明：数列 \(\{b_{n}\}\) 是等差数列；
\item 设 \(c_{n} = (a_{n + 1} - a_{n})q^{n - 1}\)（\(q\ne 0\)，\(\ n\in \mathbf{N}^{*}\)），求数列 \(\{c_n\}\) 的前 \(n\) 项和 \(S_{n}\).
\end{enumerate}
\end{problem}

\begin{solution}
\begin{enumerate}
\item 在已知递推关系中取 \((m,n)=(2,1)\)，得
\[
a_3+a_1=2a_2+2(2-1)^2=4+2=6
\]
所以 \(a_3=6\). 再取 \((m,n)=(3,1)\)，得
\[
a_5+a_1=2a_3+2(3-1)^2=12+8=20
\]
所以 \(a_5=20\).

\item 在递推关系中取 \((m,n)=(n+2,n)\)，得到
\[
a_{2n+3}+a_{2n-1}=2a_{2n+1}+2\cdot2^2=2a_{2n+1}+8
\]
移项得
\[
a_{2n+3}-a_{2n+1}=a_{2n+1}-a_{2n-1}+8
\]
按 \(b_n=a_{2n+1}-a_{2n-1}\) 的定义，这就是
\[
b_{n+1}=b_n+8
\]
又 \(b_1=a_3-a_1=6\)，所以 \(\{b_n\}\) 是首项为 \(6\)、公差为 \(8\) 的等差数列.

\item 由上一问
\[
b_n=6+8(n-1)=8n-2
\]
因此对 \(n\geq1\)，
\[
a_{2n-1}=a_1+\sum_{j=1}^{n-1}b_j
=\sum_{j=1}^{n-1}(8j-2)
=(2n-2)(2n-1)
\]
在原递推关系中取 \((m,n)=(n+1,n)\)，得
\[
a_{2n+1}+a_{2n-1}=2a_{2n}+2
\]
代入奇数项的表达式，得到
\[
a_{2n}=2n(2n-1)
\]
故对所有正整数 \(n\)，都有 \(a_n=n(n-1)\). 于是
\(a_{n+1}-a_n=2n\)，\(qquad c_n=2nq^{n-1}\).
当 \(q=1\) 时，
\[
S_n=2\sum_{k=1}^{n}k=n(n+1)
\]
当 \(q\ne1\) 时，由等比数列求和公式对 \(1+q+\cdots+q^n\) 求导，得
\[
S_n=
\begin{cases}
n(n+1), & q=1,\\
\dfrac{2\left(nq^{n+1}-(n+1)q^n+1\right)}{\left(q-1\right)^2}, & q\ne 1.
\end{cases}
\]
\end{enumerate}
\end{solution}

\begin{problem}
设 \(f(x) = \frac{1 + a^x}{1 - a^x}\)（\(a > 0\) 且 \(a \ne 1\)），\(g(x)\) 是 \(f(x)\) 的反函数.
\begin{enumerate}
\item 设关于 \(x\) 的方程 \(\log_{a} \frac{t}{(x^2 - 1)(7 - x)} = g(x)\) 在区间 \([2,6]\) 上有实数解，求 \(t\) 的取值范围.
\item 当 \(a = \e\)（\(\e\) 为自然对数的底数）时，证明：\(\sum_{k=2}^{n} g(k) > \frac{2 - n - n^2}{\sqrt{2n(n+1)}}\)；
\item 当 \(0 < a \leqslant \frac{1}{2}\) 时，试比较 \(\left| \sum_{k=1}^{n} f(k) - n \right|\) 与 \(4\) 的大小，并说明理由.
\end{enumerate}
\end{problem}

\begin{solution}
\begin{enumerate}
\item 由
\[
y=\frac{1+a^x}{1-a^x}
\]
解得 \(a^x=\frac{y-1}{y+1}\)，所以反函数为
\(g(y)=\log_a\frac{y-1}{y+1}\)，\(y\in(-\infty,-1)\cup(1,\infty)\).
在 \(x\in[2,6]\) 时，\(g(x)\) 有定义. 原方程有解等价于
\[
\frac{t}{(x^2-1)(7-x)}=a^{g(x)}=\frac{x-1}{x+1}
\]
由于 \(x\in[2,6]\) 时 \((x^2-1)(7-x)>0\)，这也保证对数的真数为正. 整理得
\[
t=(x-1)^2(7-x)
\]
令 \(h(x)=(x-1)^2(7-x)\)，则
\[
h'(x)=3(x-1)(5-x)
\]
所以 \(h\) 在 \([2,5]\) 上递增，在 \([5,6]\) 上递减，且
\(h(2)=5\)，\(qquad h(5)=32\)，\(qquad h(6)=25\).
故 \(t\) 的取值范围为 \([5,32]\).

\item 当 \(a=\e\) 时，
\[
g(k)=\ln\frac{k-1}{k+1}
\]
因此
\[
\sum_{k=2}^{n}g(k)
=\ln\prod_{k=2}^{n}\frac{k-1}{k+1}
=\ln\frac{2}{n(n+1)}
\]
这里按求和下标取 \(n\geq2\)；若取 \(n=1\)，题中的严格不等式退化为等式. 令
\[
u=\frac{n(n+1)}2>1
\]
则
\(\ln\frac{2}{n(n+1)}=-\ln u\)，\(\frac{2-n-n^2}{\sqrt{2n(n+1)}}=-\frac{u-1}{\sqrt u}\).
设
\[
\phi(u)=\sqrt u-\frac1{\sqrt u}-\ln u
\]
则 \(\phi(1)=0\)，且
\(\phi'(u)=\frac{(\sqrt u-1)^2}{2u^{3/2}}>0\)，\((u>1)\).
所以 \(\ln u<\sqrt u-\frac1{\sqrt u}\)，即
\[
\sum_{k=2}^{n}g(k)>\frac{2-n-n^2}{\sqrt{2n(n+1)}}
\]

\item 因为 \(0<a\leq\frac12\)，对每个正整数 \(k\) 有 \(0<a^k<1\)，从而
\[
f(k)-1=\frac{2a^k}{1-a^k}>0
\]
故
\[
\left|\sum_{k=1}^{n}f(k)-n\right|
=2\sum_{k=1}^{n}\frac{a^k}{1-a^k}
\]
又因 \(a^k\leq a\)，有
\[
\frac{a^k}{1-a^k}\leq\frac{a^k}{1-a}
\]
于是
\[
2\sum_{k=1}^{n}\frac{a^k}{1-a^k}
\leq\frac{2}{1-a}\sum_{k=1}^{n}a^k
<\frac{2}{1-a}\cdot\frac{a}{1-a}
=\frac{2a}{(1-a)^2}\leq4
\]
最后一个不等式在 \(0<a\leq\frac12\) 时成立，而前面的有限项和严格小于无穷等比级数之和，所以
\[
\left|\sum_{k=1}^{n}f(k)-n\right|<4
\]
\end{enumerate}
\end{solution}
